\newpage
\section{Specyfikacja wymagań}
\label{sec:specyfikacja}

Specyfikacja opisuje, czego oczekujemy od systemu, bez przesądzania o~sposobie
jego realizacji. Wymagania funkcjonalne oznaczono identyfikatorami WF-xx,
a~niefunkcjonalne -- WNF-xx. Każde wymaganie ma przypisany priorytet:

\begin{itemize}
    \item \textbf{MUST} -- wymaganie konieczne; bez niego system nie realizuje
    swojego podstawowego celu,
    \item \textbf{SHOULD} -- wymaganie ważne, które powinno zostać zrealizowane,
    ale system może bez niego funkcjonować,
    \item \textbf{COULD} -- wymaganie dodatkowe, realizowane w~miarę dostępnego
    czasu.
\end{itemize}

\subsection{Zakres systemu}

System obejmuje:

\begin{itemize}
    \item wyszukiwanie środków ochrony roślin po nazwie handlowej, substancji
    czynnej, uprawie i~agrofagu, wraz z~filtrowaniem krzyżowym,
    \item kartę środka z~informacjami o~zastosowaniach, dawkach, terminach
    prawnych i~etykiecie,
    \item wyszukiwanie środków o~tych samych substancjach czynnych
    (zamienników),
    \item automatyczną aktualizację danych na podstawie publicznych danych
    Ministerstwa Rolnictwa i~Rozwoju Wsi.
\end{itemize}

Poza zakresem systemu są: ceny i~sprzedaż środków, konta użytkowników i~panel
administracyjny, ewidencja wykonanych zabiegów, rozpoznawanie agrofagów na
podstawie zdjęć, aplikacja instalowana na urządzeniu oraz praca bez dostępu do
internetu.

\subsection{Interesariusze}

\begingroup
\small
\begin{longtable}{|>{\raggedright\arraybackslash}p{3.2cm}|p{8.8cm}|>{\raggedright\arraybackslash}p{2.4cm}|}
\hline
\textbf{Interesariusz} & \textbf{Opis i~potrzeba} & \textbf{Wymagania} \\
\hline
\endfirsthead
\hline
\textbf{Interesariusz} & \textbf{Opis i~potrzeba} & \textbf{Wymagania} \\
\hline
\endhead
Doradca agronomiczny / rolnik profesjonalny &
Analizuje składy chemiczne środków, planuje rotację preparatów i~weryfikuje
ich zgodność z~aktualnym prawem oraz zasadami integrowanej ochrony roślin.
,,Chcę wiedzieć, czym mogę zastąpić środek i~czy jest on dopuszczony na moją
uprawę.'' &
WF-03, WF-04, WF-10, WF-16 \\
\hline
Rolnik indywidualny / operator maszyn &
Wykonuje zabiegi w~terenie i~sprawdza na smartfonie dawkowanie, warunki
stosowania oraz okresy karencji i~prewencji bezpośrednio z~kabiny ciągnika.
,,Chcę w~kabinie ciągnika sprawdzić dawkę i~karencję, nawet przy słabym
zasięgu.'' &
WF-06, WF-07, WNF-01 do WNF-05 \\
\hline
Ogrodnik amator / działkowiec &
Szuka bezpiecznych środków do upraw przydomowych, dostępnych bez uprawnień
profesjonalnych. ,,Chcę wiedzieć, których środków mogę legalnie używać
i~czy istnieją tańsze odpowiedniki drogiej, małej butelki ze sklepu
ogrodniczego.'' &
WF-05, WF-10 \\
\hline
Ministerstwo Rolnictwa i~Rozwoju Wsi &
Dostawca danych -- publikuje rejestr środków ochrony roślin na portalu
dane.gov.pl. Jego dane muszą być prezentowane bez zniekształceń, a~rytm ich
publikacji wyznacza rytm aktualizacji systemu. &
WF-13, WNF-07, O-01 \\
\hline
Zespół projektowy &
Twórcy i~opiekunowie systemu. ,,Chcemy, żeby dane aktualizowały się bez
naszej ręcznej pracy.'' &
WF-13, WNF-08, O-04 \\
\hline
Producenci i~dystrybutorzy środków (interesariusz pośredni) &
Ich produkty są prezentowane w~systemie, dlatego oczekują rzetelnej
i~neutralnej prezentacji. &
WNF-09 \\
\hline
\end{longtable}
\endgroup

\subsection{Role użytkowników}

System nie wymaga zakładania kont. Użytkownik wybiera swoją rolę przełącznikiem
w~aplikacji, a~wybór jest zapamiętywany na jego urządzeniu. System nie
weryfikuje roli -- ma ona charakter informacyjny i~decyduje o~tym, jakie środki
są prezentowane.

\begin{itemize}
    \item \textbf{Użytkownik profesjonalny} (doradca agronomiczny, rolnik,
    operator maszyn) -- ma dostęp do wszystkich środków zarejestrowanych
    w~Polsce.
    \item \textbf{Użytkownik nieprofesjonalny} (ogrodnik amator, działkowiec)
    -- widzi wyłącznie środki dopuszczone do stosowania przez użytkowników
    nieprofesjonalnych, również na liście zamienników.
\end{itemize}

\subsection{Wymagania funkcjonalne}

\begingroup
\small
\begin{longtable}{|p{1.4cm}|p{11.4cm}|p{1.6cm}|}
\hline
\textbf{ID} & \textbf{Wymaganie} & \textbf{Priorytet} \\
\hline
\endfirsthead
\hline
\textbf{ID} & \textbf{Wymaganie} & \textbf{Priorytet} \\
\hline
\endhead
WF-01 & Użytkownik może wyszukać środek po nazwie handlowej. Wyniki są
wyświetlane po wpisaniu co najmniej 3 znaków. & MUST \\
\hline
WF-02 & System odnajduje środek, uprawę, agrofaga lub substancję czynną,
gdy wpisana nazwa różni się od nazwy w~rejestrze najwyżej o~2 znaki
(brakujący, nadmiarowy, błędny lub zamieniony znak), niezależnie od
wielkości liter i~polskich znaków. & MUST \\
\hline
WF-03 & Użytkownik może wyszukać środki zawierające wybraną substancję
czynną. & MUST \\
\hline
WF-04 & Użytkownik może jednocześnie wybrać uprawę, agrofaga i~rodzaj środka.
Wyniki zawierają wyłącznie środki, które mają zastosowanie spełniające
wszystkie wybrane kryteria, a~każdy środek występuje na liście jeden raz.
& MUST \\
\hline
WF-05 & Użytkownik może wybrać rolę profesjonalną lub nieprofesjonalną.
W~roli nieprofesjonalnej wyniki wyszukiwania i~lista zamienników obejmują
wyłącznie środki dopuszczone do stosowania przez użytkowników
nieprofesjonalnych. & MUST \\
\hline
WF-06 & Karta środka zawiera: nazwę, producenta, numer zezwolenia, rodzaj
środka, substancje czynne wraz z~zawartością, termin ważności zezwolenia,
terminy dopuszczenia do sprzedaży i~stosowania, dostępne opakowania oraz
listę wszystkich zastosowań (uprawa, agrofag, dawka, termin i~metoda
stosowania). & MUST \\
\hline
WF-07 & Karta środka wyświetla okres karencji, jeżeli jest on podany
w~danych źródłowych. W~przeciwnym razie informuje, że okres karencji należy
sprawdzić na etykiecie, i~udostępnia odnośnik do niej. & MUST \\
\hline
WF-08 & System odczytuje okres karencji z~etykiet środków i~wyświetla go na
karcie środka dla każdej uprawy. & COULD \\
\hline
WF-09 & Użytkownik może wyświetlić etykietę środka bez opuszczania aplikacji
oraz pobrać ją na swoje urządzenie. & SHOULD \\
\hline
WF-10 & Dla wybranego środka system wyświetla inne środki zawierające te
same substancje czynne (zamienniki), podając przy każdym z~nich zawartość
substancji czynnej, dostępne opakowania oraz rodzaj użytkownika
(profesjonalny lub nieprofesjonalny). Środki o~tej samej zawartości
substancji czynnej są wyświetlane na początku listy, a~po wybraniu uprawy
lista obejmuje wyłącznie środki zarejestrowane dla tej uprawy. & MUST \\
\hline
WF-11 & System oznacza środki, których termin dopuszczenia do stosowania
upływa w~ciągu 90 dni lub już upłynął. & SHOULD \\
\hline
WF-12 & Wyszukanie środka usuniętego z~rejestru zwraca informację o~jego
usunięciu wraz z~datą, zamiast pustej listy wyników. & SHOULD \\
\hline
WF-13 & System automatycznie pobiera i~wczytuje nowe wersje danych
o~środkach ochrony roślin publikowane przez Ministerstwo Rolnictwa i~Rozwoju
Wsi na portalu dane.gov.pl. & MUST \\
\hline
WF-14 & System wyświetla datę publikacji wersji danych, z~której korzysta.
& SHOULD \\
\hline
WF-15 & System zapamiętuje na urządzeniu użytkownika 10 ostatnio
przeglądanych środków, bez zakładania konta. & COULD \\
\hline
WF-16 & Każda karta środka ma stały adres, który użytkownik może przesłać
innej osobie. & COULD \\
\hline
\end{longtable}
\endgroup

\subsection{Wymagania niefunkcjonalne}

\begingroup
\small
\begin{longtable}{|p{1.4cm}|p{11.4cm}|p{1.6cm}|}
\hline
\textbf{ID} & \textbf{Wymaganie} & \textbf{Priorytet} \\
\hline
\endfirsthead
\hline
\textbf{ID} & \textbf{Wymaganie} & \textbf{Priorytet} \\
\hline
\endhead
WNF-01 & Wyniki wyszukiwania są wyświetlane w~czasie nie dłuższym niż
2~sekundy od zakończenia wpisywania zapytania, przy połączeniu
o~przepustowości 1~Mb/s. & MUST \\
\hline
WNF-02 & Pierwsze otwarcie aplikacji wymaga pobrania nie więcej niż 1~MB
danych (dla porównania: wyszukiwarka ministerstwa -- ok.~10~MB). & MUST \\
\hline
WNF-03 & Żaden widok aplikacji nie wymaga przewijania w~poziomie na ekranach
o~szerokości od 360 pikseli. & MUST \\
\hline
WNF-04 & Elementy dotykowe (przyciski, pola, pozycje list) mają wymiary co
najmniej $48 \times 48$ pikseli, a~odstęp między nimi wynosi co najmniej
8 pikseli. & SHOULD \\
\hline
WNF-05 & Kontrast tekstu względem tła wynosi co najmniej 7:1 (poziom AAA
według wytycznych WCAG~2.1). & SHOULD \\
\hline
WNF-06 & Aplikacja działa poprawnie w~dwóch najnowszych wersjach przeglądarek
Chrome, Firefox, Safari i~Edge, na systemach Android, iOS i~Windows. & MUST \\
\hline
WNF-07 & Nowa wersja danych jest dostępna w~systemie najpóźniej 7~dni po jej
publikacji na portalu dane.gov.pl. & MUST \\
\hline
WNF-08 & Nieudane wczytanie nowej wersji danych nie przerywa działania
systemu: użytkownicy nadal korzystają z~poprzedniej wersji danych, a~informacja
o~błędzie jest zapisywana w~dzienniku zdarzeń systemu. & MUST \\
\hline
WNF-09 & Wszystkie funkcje systemu są dostępne bez zakładania konta
i~logowania, bez opłat oraz bez wyświetlania reklam. & MUST \\
\hline
WNF-10 & Interfejs używa terminologii zgodnej z~rejestrem ministerstwa
i~etykietami środków (np. ,,substancja czynna'', ,,dawka'', ,,okres
karencji''). & SHOULD \\
\hline
\end{longtable}
\endgroup

\subsection{Ograniczenia}

\begin{itemize}
    \item \textbf{O-01:} Jedynym źródłem danych są pliki publikowane przez
    Ministerstwo Rolnictwa i~Rozwoju Wsi na portalu dane.gov.pl. System nie
    zawiera informacji spoza tych plików, np. cen środków, a~okres karencji
    jest w~nich podany tylko w~części opisów zastosowań.
    \item \textbf{O-02:} System ma charakter wyłącznie informacyjny -- każda
    karta środka zawiera informację, że prawnie wiążąca jest
    etykieta-instrukcja stosowania środka.
    \item \textbf{O-03:} System jest aplikacją internetową uruchamianą
    w~przeglądarce; nie przewiduje się aplikacji instalowanej na urządzeniu.
    \item \textbf{O-04:} System nie ma panelu administracyjnego, a~aktualizacja
    danych odbywa się w~pełni automatycznie.
    \item \textbf{O-05:} Projekt realizuje zespół dwuosobowy w~ciągu jednego
    semestru, dlatego wymagania o~priorytecie COULD są realizowane w~miarę
    dostępnego czasu.
\end{itemize}

\subsection{Kryteria akceptacji}

\subsubsection{WF-02 -- wyszukiwanie odporne na literówki}

\begin{enumerate}
    \item Wpisanie ,,Amistr'', ,,Amitsar'' lub ,,Amistarr'' zwraca środek
    Amistar 250 SC wśród pierwszych 5 wyników.
    \item Zapytania ,,AMISTAR'', ,,amistar'' i~,,Amistar'' zwracają te same
    wyniki.
    \item Wpisanie ,,pszenzyto'' w~polu uprawy odnajduje uprawę ,,pszenżyto''.
    \item Gdy żaden środek nie pasuje do zapytania, system wyświetla komunikat
    o~braku wyników.
\end{enumerate}

\subsubsection{WF-04 -- filtrowanie krzyżowe}

\begin{enumerate}
    \item Wybranie uprawy ,,ziemniak'' i~agrofaga ,,stonka ziemniaczana''
    zwraca wyłącznie środki, które w~rejestrze mają zastosowanie dla tej
    pary.
    \item Każdy środek pojawia się na liście wyników jeden raz, niezależnie od
    liczby jego zastosowań.
    \item Dodanie filtra rodzaju środka ,,Insektycyd'' zawęża listę -- każdy
    wyświetlony środek ma rodzaj obejmujący ,,Insektycyd''.
    \item Jeżeli żaden środek nie spełnia wszystkich kryteriów, system
    wyświetla komunikat o~braku wyników i~umożliwia usunięcie pojedynczego
    filtra jednym dotknięciem.
\end{enumerate}

\subsubsection{WF-10 -- zamienniki}

\begin{enumerate}
    \item Dla środka Amistar 250 SC (azoksystrobina -- 250~g) lista zawiera
    wyłącznie środki, których jedyną substancją czynną jest azoksystrobina,
    a~środki o~zawartości 250~g znajdują się na jej początku.
    \item Przy każdym środku na liście wyświetlana jest zawartość substancji
    czynnej, dostępne opakowania oraz rodzaj użytkownika.
    \item W~roli nieprofesjonalnej lista zawiera wyłącznie środki dopuszczone
    dla użytkowników nieprofesjonalnych; jeżeli takich środków nie ma, system
    wyświetla komunikat o~braku zamienników dla użytkowników
    nieprofesjonalnych.
    \item Po wybraniu uprawy lista zawiera wyłącznie środki zarejestrowane dla
    tej uprawy.
    \item Środek, dla którego wyszukiwane są zamienniki, nie pojawia się na
    liście, a~środki, których termin dopuszczenia do stosowania upłynął, nie są
    proponowane.
\end{enumerate}

\subsection{Konflikty wymagań}

\begin{itemize}
    \item \textbf{K-01: tańsze odpowiedniki a~przepisy.} Problem opisany
    w~sekcji~\ref{sec:problem} zakłada, że użytkownik nieprofesjonalny
    znajdzie tańszy odpowiednik środka w~większym opakowaniu.
    Według danych ministerstwa tylko około 7\,\% środków w~rejestrze jest
    dopuszczonych do stosowania przez użytkowników nieprofesjonalnych --
    pozostałe wymagają uprawnień, niezależnie od wielkości opakowania. Rozstrzygnięcie: przy każdym
    zamienniku wyświetlany jest rodzaj użytkownika (WF-10), a~w~roli
    nieprofesjonalnej prezentowane są tylko środki dopuszczone dla amatorów
    (WF-05), dzięki czemu użytkownik widzi wyłącznie te odpowiedniki, które
    może legalnie kupić i~stosować.
\end{itemize}
